\chapter{Water accounting and time stepping}\label{ch:c5}
This chapter explains how the two models keep their water accounts consistent, what exactly happens in each time step,
and how runs are started, restarted and repeated.

\begin{plain}
Every day the two models take turns: Raven works out how much water goes down; MODFLOW moves it through the aquifer and
reports what comes back up; Raven then books the returned water. Both keep a running account, and the two accounts are
compared every day. This chapter also explains how to stop a run and carry on later (a restart), and how to run many
versions of a model at once (an ensemble).
\end{plain}

\section{Keeping the books: water balance}\label{ch:balance}
\subsection{MODFLOW's balance}
For every step Raven computes all boundary flows from MODFLOW's own matrix terms ($q=\mathrm{HCOF}\cdot h-\mathrm{RHS}$)
and the storage change from MODFLOW's storage package, and reports the imbalance
\begin{multline}
\varepsilon=Q^{\mathrm{rch}}+Q^{\mathrm{sto}}-Q^{\mathrm{seep}}+Q^{\mathrm{riv}}_{\mathrm{in}}-Q^{\mathrm{riv}}_{\mathrm{out}}-Q^{\mathrm{evt}}
+Q^{\mathrm{wel}}+Q^{\mathrm{ghb}}+Q^{\mathrm{chd}}\\+Q^{\mathrm{swx}}_{\mathrm{in}}-Q^{\mathrm{swx}}_{\mathrm{out}}+Q^{\mathrm{res}},
\end{multline}
where $Q^{\mathrm{sto}}$ is water released from storage. It is written in m$^3$/d and as a percentage of the gross flow
(the sum of the absolute values of all terms). In the tests of Chapter~\ref{ch:verification} the worst single-day error is about $5\times10^{-6}$\% and the
cumulative error about $10^{-10}$.

\subsection{Raven's balance}
Raven tracks cumulative input $I$ and output $O$ and reports $\mathrm{MB}=\left(S-S_0\right)+\left(O-I\right)$ in mm over
the watershed. The coupling adds:
\begin{itemize}
\item \textbf{to output}: recharge (the end-of-day \cmd{GROUNDWATER} value of coupled HRUs, as Raven already does), and
reservoir seepage (already counted by Raven as a reservoir loss);
\item \textbf{to input}: every volume the coupling adds to Raven in the step --- seepage returned to HRU stores, the net
river exchange applied to reaches, and the net surface-store exchange. A net loss (for example a losing river) is a
negative input.
\end{itemize}
With these, Raven's own balance closes exactly as without groundwater (0.43--0.52 mm maximum over 20 years in the Liard
example, the same as uncoupled).

\subsection{The whole system}
The water that leaves Raven as recharge enters MODFLOW (directly or via the delay store); the water MODFLOW sends back
enters Raven. Reservoir seepage enters MODFLOW one day after leaving Raven. Therefore
\begin{equation}
\sum V^{\mathrm{Raven}}_{\mathrm{rch}}=\sum V^{\mathrm{MF6}}_{\mathrm{rch}}+\Delta S^{\mathrm{delay}}+\sum V^{\mathrm{CHD}}_{\mathrm{rch}},
\end{equation}
and the audit tool checks this identity, the river, seepage and store bookkeeping, and MODFLOW's balance for every run.

\section{A real day of water accounts}
\begin{example}[: 15 June 1996 in the Liard run]
From \cmd{GWBudget.csv} (all in m$^3$/d):
\begin{center}\begin{tabular}{lr}\toprule\hdr
Inflows to the aquifer & \\\midrule
recharge from Raven & 21\,365\,037\\
river leakage & 2\,978\,677\\\midrule\hdr
Outflows from the aquifer & \\\midrule
seepage to Raven & 19\,087\,047\\
discharge to rivers & 1\,331\,456\\\midrule
storage release (negative = stored) & $-$3\,925\,212\\\bottomrule
\end{tabular}\end{center}
Inflows $-$ outflows $+$ release $=24\,343\,714-20\,418\,502-3\,925\,212=0.00003$ m$^3$/d, i.e.\ a relative error of
about $10^{-10}$. Raven recorded the same numbers on its side: 21.4 million m$^3$ of recharge leaving its stores, 19.1 million m$^3$
of seepage entering them, and a net river exchange of $-1.65$ million m$^3$ applied to its reaches (the rivers lost more than
they gained that day). The aquifer stored the difference: a water-table rise of about 1 mm on average.
\end{example}

\section{How to read the balance outputs}
\begin{enumerate}
\item \emph{MF6 balance error [\%]} should stay well below 0.1\% every step; larger values point to a non-converged step or a
too-loose head tolerance (\cmd{:SolverHead\allowbreak Tolerance}, default 0.00001 m).
\item \emph{Raven recharge} and \emph{MF6 recharge} are equal unless a recharge delay or fixed-head cells are used.
\item \emph{returned to Raven} equals \emph{seepage to Raven}.
\item \emph{river loss carried to next step} is normally zero; when not, it is taken from the river the next day.
\item Raven's \cmd{Watershed\allowbreak Storage.csv} MB error should be as small as without groundwater.
\end{enumerate}
The tool \cmd{tools\_audit.py} checks all of these automatically.

\section{The coupled time step, start-up and restarts}\label{ch:timestep}
\subsection{The algorithm}
\begin{figure}[!htbp]\centering
\begin{tikzpicture}[x=1cm,y=1cm,b/.style={draw,rounded corners,font=\scriptsize,align=center,minimum height=8mm,text width=23mm}]
\node[b,fill=teal!12] (a) at (0,0) {Raven HRU\\processes};
\node[b,fill=gray!15] (b) at (2.9,0) {recharge and\\stresses to MF6};
\node[b,fill=violet!12] (c) at (5.8,0) {MODFLOW\\Newton solve};
\node[b,fill=gray!15] (d) at (8.7,0) {seepage, stores,\\river terms back};
\node[b,fill=teal!12] (e) at (11.6,0) {Raven routing\\and balance};
\foreach \u/\v in {a/b,b/c,c/d,d/e} \draw[-{Stealth[length=2mm]},thick] (\u) -- (\v);
\draw[->,thick,gray] (-1.2,-0.8) -- (12.8,-0.8) node[right,font=\scriptsize]{time};
\end{tikzpicture}
\caption{Order of work within one time step.}\label{fig:timeline}
\end{figure}

Algorithm~\ref{alg:step} lists exactly what happens in one Raven time step of a coupled model.

\begin{algorithm}[h]\caption{One coupled time step}\label{alg:step}
\begin{algorithmic}[1]
\State Raven: vertical and lateral processes of all HRUs; coupled HRUs' \cmd{GROUNDWATER} now holds the day's recharge
\State Complete yesterday's line of \cmd{GWBudget.csv} (river and reservoir terms are final only after routing)
\State Recharge: delay reservoirs (Ch.~\ref{ch:recharge}); spread over cells; skip fixed-head cells
\State MODFLOW \cmd{prepare\_time\_step}; set RCH rates
\State Set river stages and conductances, well rates, GHB/CHD heads (Ch.~\ref{ch:rivers},\ref{ch:wells})
\State Set water-table ET rates from PET left unmet by Raven (Ch.~\ref{ch:et})
\State Deliver yesterday's reservoir seepage (Ch.~\ref{ch:reservoirs}); set surface-store levels and limited conductances (Ch.~\ref{ch:stores})
\State MODFLOW \cmd{prepare\_solve}; \textbf{repeat} \cmd{solve} \textbf{until} converged or the iteration limit; \cmd{finalize\_solve}; \cmd{finalize\_time\_step}
\State Compute every package flow as $\mathrm{HCOF}\cdot h-\mathrm{RHS}$; reservoir local heads for today's routing
\State Surface-store exchange applied to Raven stores; seepage and river flows computed
\If{the step did not converge} limit seepage and river gains returned to Raven to the step's total inflow; warn \EndIf
\State Seepage added to Raven stores; balance, heads, connectivity and NetCDF output written
\State Raven: before routing, river gains added to / losses taken from each subbasin's in-catchment inflow
\State Raven routing; losses not covered are taken from the flow entering each reach; reservoirs compute seepage
\State Raven: water balance (recharge out, all returned water in); diagnostics including \cmd{GW\_HEAD}
\end{algorithmic}
\end{algorithm}

\subsection{Starting conditions}\label{sec:init}
Initial heads are chosen in this order: (1) a hotstart block \cmd{:GWHeads} in the \cmd{.rvc} file; (2)
\cmd{:InitialGWHeads} in the \cmd{.rvc} file (\cmd{DEPTH\_BELOW\_SURFACE d} or \cmd{ELEVATION z}); (3) each profile's
\cmd{INITIAL\_HEAD\_DEPTH}. Guessed initial heads can be far from reality; \cmd{:GWInitialization STEADY\_STATE r}
first solves a steady state before Raven's first day, with recharge $r$ [mm/d] over the area covered by coupled HRUs
(cell $i$ receives $r\,a_i/1000$ m$^3$/d) and the initial river stages, well rates and boundary heads. This is skipped on a hotstart. A steady state may not
exist --- for example when wells pump more than the aquifer can supply --- and then the solve does not converge; Raven
warns and starts from the initial heads instead, as if \cmd{:GWInitialization} had not been given. The groundwater model is built after \emph{all} input files,
including the \cmd{.rvc}, have been read.

\subsection{Restarting a run (hotstart)}
At the end of a run Raven writes \cmd{solution.rvc}. Besides Raven's own states it contains \cmd{:GWHeads} (the head in
every cell), \cmd{:GWRechargeStore} (recharge delay stores, if used), \cmd{:GWRiverLossCarry} (river loss carried, if any) and \cmd{:GWReservoirSeepage} (reservoir seepage in
transit, if used). Using this file as the \cmd{.rvc} of a new run --- with \cmd{:StartDate} set to the end of the first
run --- continues exactly: in the tests the restarted run matched a continuous run to within $6\times10^{-6}$ in all flows.

\subsection{Ensembles}
Raven's ensemble modes (Monte Carlo, DDS, EnKF) run the model several times in one program run. For every member the
groundwater model is rebuilt from scratch, MODFLOW is restarted, and its files go to that member's output folder.
Raven itself was corrected so that each member starts with a fresh water balance and, unless its own \cmd{.rvc} says
otherwise, from the routing memory and reservoir states of the first run (stage, outflow, loss terms). These defaults are
restored before the member's \cmd{.rvc} is read, so member-specific initial conditions (as used by EnKF forecasting) still
apply; groundwater hotstart values are likewise cleared, so each member uses only its own \cmd{.rvc}, and the steady-state
start is decided for each member separately
(Section~\ref{sec:coreflix}).

\subsection{Safeguards}\label{sec:guard}
\begin{itemize}
\item \textbf{Non-converged steps.} All water returned to Raven in such a step (seepage plus river gains) is limited to
the water that entered the aquifer in that step (recharge, river and store leakage, reservoir injection and boundary
inflows); discharge into wetland and lake stores is limited the same way. The step is flagged and a warning written.
\item \textbf{Rivers and stores} can never lose more water than they hold (Sections~\ref{ch:rivers},\ref{ch:stores}).
\item \textbf{Wells} cannot pump from a dry cell (\cmd{AUTO\_FLOW\_REDUCE}).
\item \textbf{Step count}: Raven stops if it would run past MODFLOW's last step.
\item \textbf{MODFLOW folder}: if Raven cannot enter the MODFLOW folder, the call fails with a clear message instead of
running MODFLOW in the wrong place.
\end{itemize}
